Section \ref{sec.stability} describes the process in the mean.
This section computes its stationary second moments,
in closed form, for the exponential kernel.
Throughout, $\branchingRatio < 1$ and $\baseIntensity > 0$,
so that Proposition \ref{prop.stationarity} applies
and the process is taken in its stationary version.

\subsubsection{The covariance of the state}

Theorem \ref{thm.markovState} is applied with $f(z) = z z^{\transpose}$.
At an event of type $e$ the state gains the unit vector $\unitVector_e$,
so $\decayedCounts\decayedCounts^{\transpose}$ gains
$\decayedCounts \unitVector_e^{\transpose}
+ \unitVector_e \decayedCounts^{\transpose}
+ \unitVector_e \unitVector_e^{\transpose}$.
Two properties of the jumps enter here, and both are hypotheses rather than conveniences:
the jumps of $\decayedCounts$ are \emph{unit} vectors,
which is what makes the last term $\unitVector_e \unitVector_e^{\transpose}$;
and no two components jump together, by \eqref{eq.noCommonJumps},
which is what keeps that term diagonal.

\begin{thm}\label{thm.lyapunov}
 Let $\excitation \geq 0$, $\decay > 0$, $\baseIntensity > 0$ and $\branchingRatio < 1$,
 and write $\hurwitzMatrix := \excitation - \decay I$.
 The stationary version of Proposition \ref{prop.stationarity} has finite second moments,
 its covariance $V := \CoVariance(\decayedCounts)$ is the unique solution of the Lyapunov
 equation
 \begin{equation}\label{eq.lyapunov}
  \hurwitzMatrix V + V \hurwitzMatrix^{\transpose} + \diag(\stationaryIntensity) = 0 ,
 \end{equation}
 it is given by
 $V = \int_{0}^{\infty} e^{\hurwitzMatrix s} \diag(\stationaryIntensity)
 e^{\hurwitzMatrix^{\transpose} s} \, ds$ and is positive definite,
 and the stationary covariance of the intensity is
 $\CoVariance(\intensity[ ]) = \excitation V \excitation^{\transpose}$.
\end{thm}
\begin{proof}
 Start the process at $\decayedCounts(0) = \multiCountingProc(0) = 0$ and write
 $m_t := \Expectation[\decayedCounts(t)]$ and
 $\secondMoment_t := \Expectation[\decayedCounts(t)\decayedCounts(t)^{\transpose}]$,
 both finite for every $t$ by Remark \ref{remark.noSubcriticality}.
 Theorem \ref{thm.markovState} with $f(z)=z$ and \eqref{eq.intensityFromState} give
 $\dot m_t = \hurwitzMatrix m_t + \baseIntensity$,
 and with $f(z) = zz^{\transpose}$, using
 $\Expectation[\decayedCounts \intensity[ ]^{\transpose}]
 = m_t\baseIntensity^{\transpose} + \secondMoment_t\excitation^{\transpose}$,
 \begin{equation}\label{eq.secondMomentODE}
  \dot{\secondMoment}_t
  = \hurwitzMatrix \secondMoment_t + \secondMoment_t \hurwitzMatrix^{\transpose}
  + \baseIntensity m_t^{\transpose} + m_t \baseIntensity^{\transpose}
  + \diag\big( \baseIntensity + \excitation m_t \big) .
 \end{equation}
 By Proposition \ref{prop.lyapunovUniqueness}~\emph{iii.},
 $\branchingRatio<1$ makes $\hurwitzMatrix$ Hurwitz,
 so $m_t \to m = -\hurwitzMatrix^{-1}\baseIntensity = \stationaryIntensity/\decay$
 and the linear equation \eqref{eq.secondMomentODE} converges to its unique fixed point
 $\secondMoment$, which by Proposition \ref{prop.lyapunovUniqueness}~\emph{i.--ii.}
 exists and is unique.
 A stationary version exists by Proposition \ref{prop.stationarity},
 and its second moment is that fixed point, the limit being independent of the initial state.
 Writing $\secondMoment = V + m m^{\transpose}$ and using
 $\hurwitzMatrix m = -\baseIntensity$,
 \begin{equation*}
  \hurwitzMatrix m m^{\transpose} + m m^{\transpose}\hurwitzMatrix^{\transpose}
  = -\baseIntensity m^{\transpose} - m\baseIntensity^{\transpose} ,
 \end{equation*}
 so the two cross terms of \eqref{eq.secondMomentODE} cancel and what remains is
 \eqref{eq.lyapunov}, with $\diag(\baseIntensity + \excitation m)
 = \diag(\stationaryIntensity)$ by \eqref{eq.intensityFromState}.
 The integral formula and positive definiteness are
 Proposition \ref{prop.lyapunovUniqueness}~\emph{ii.},
 the forcing term being diagonal with strictly positive entries.
 Finally $\intensity[ ] = \baseIntensity + \excitation\decayedCounts$ with
 $\baseIntensity$ constant, so
 $\CoVariance(\intensity[ ]) = \excitation V \excitation^{\transpose}$.
\end{proof}

\begin{remark}
 Notice that the forcing term of \eqref{eq.lyapunov} is the diagonal matrix of stationary
 intensities:
 fluctuation is injected at the rate at which events occur,
 one unit of variance per event,
 and with no covariance between types because no two of them fire together.
 The homogeneous part transports it,
 and the equation is the balance between the two.
\end{remark}

\begin{prop}\label{prop.lyapunovUniqueness}
 Let $\hurwitzMatrix$ be a real $d \times d$ matrix.
 A matrix is called \emph{Hurwitz} if all its eigenvalues have strictly negative real part.
 Then
 \begin{enumerate}[label={\emph{\roman{*}}.}]
  \item the map $V \mapsto \hurwitzMatrix V + V\hurwitzMatrix^{\transpose}$ is a bijection
        on symmetric matrices if and only if no two eigenvalues of $\hurwitzMatrix$ sum to
        zero; in particular a Hurwitz $\hurwitzMatrix$ is sufficient, and not necessary,
        for $\hurwitzMatrix V + V\hurwitzMatrix^{\transpose} + Q = 0$ to have a unique
        solution;
  \item if $\hurwitzMatrix$ is Hurwitz, that unique solution is
        $V = \int_{0}^{\infty} e^{\hurwitzMatrix s} Q e^{\hurwitzMatrix^{\transpose} s} ds$,
        and $Q \succ 0$ implies $V \succ 0$;
  \item if moreover $\excitation \geq 0$ and $\decay > 0$, then
        $\hurwitzMatrix = \excitation - \decay I$ is Hurwitz if and only if
        $\branchingRatio < 1$.
 \end{enumerate}
\end{prop}
\begin{proof}
 The eigenvalues of the map in \emph{i.} are the sums $\theta_i + \theta_j$ of the
 eigenvalues of $\hurwitzMatrix$, which gives the stated criterion;
 see \cite[Theorem 2.4.4.1]{HJ13mat}.
 For \emph{ii.}, the integral converges because $\hurwitzMatrix$ is Hurwitz,
 and differentiating $e^{\hurwitzMatrix s}Qe^{\hurwitzMatrix^{\transpose}s}$ and
 integrating over $[0,\infty)$ gives $-Q$ on the left-hand side of the equation.
 For \emph{iii.}, the eigenvalues of $\hurwitzMatrix$ are $\decay(\gamma_i - 1)$ over the
 eigenvalues $\gamma_i$ of $\branchingMatrix$,
 so $\hurwitzMatrix$ is Hurwitz if and only if $\Real \gamma_i < 1$ for every $i$.
 That is weaker than $\branchingRatio<1$ for a general matrix;
 here $\branchingMatrix \geq 0$, so by Lemma \ref{lemma.perronFrobenius}~\emph{i.}
 the spectral radius $\branchingRatio$ is itself an eigenvalue,
 and $\Real \gamma_i \leq \vert \gamma_i \vert \leq \branchingRatio$ for every $i$.
 The two conditions therefore coincide.
\end{proof}

\begin{remark}
 Notice that $\excitation \geq 0$ is what makes \emph{iii.} an equivalence.
 Without it the implication runs one way only:
 $\branchingMatrix$ with eigenvalues $\pm 2i$ has $\branchingRatio = 2$
 while $\Real\gamma_i = 0 < 1$.
\end{remark}

\subsubsection{The scalar case}

Contracting \eqref{eq.lyapunov} with $\mathbf{1}^{\transpose}$ on the left and
$\mathbf{1}$ on the right gives a scalar equation as soon as
$\mathbf{1}^{\transpose}\excitation$ is proportional to $\mathbf{1}^{\transpose}$,
that is, as soon as $\excitation$ has constant column sums.
Those sums are then necessarily $\decay\branchingRatio$:
by Lemma \ref{lemma.perronFrobenius}~\emph{iii.} applied to $\branchingMatrix^{\transpose}$,
constant column sums force $\branchingRatio$ to equal them.
Hence
$\mathbf{1}^{\transpose}\hurwitzMatrix
= -\decay(1-\branchingRatio)\mathbf{1}^{\transpose}$, and
\begin{equation}\label{eq.scalarVariance}
 \Variance \big( \mathbf{1}^{\transpose}\decayedCounts \big)
 = \frac{\totalRate}{2\decay(1-\branchingRatio)} .
\end{equation}
In one dimension this is
$\Variance(\decayedCounts) = \stationaryIntensity/(2\decay(1-\branchingRatio))$,
and it is the formula usually quoted.

\begin{remark}\label{remark.scalarVariance}
 Notice what \eqref{eq.scalarVariance} is a statement about.
 It concerns $\mathbf{1}^{\transpose}\decayedCounts$, the total decayed count,
 and not any single component;
 it is not a statement about the total intensity,
 which in the same special case has variance
 $\decay\branchingRatio^{2}\totalRate/(2(1-\branchingRatio))$;
 and it requires constant column sums, which are not generic.
 Neither parametrisation of Section \ref{sec.orderFlowModel} has them,
 and on the two of them \eqref{eq.scalarVariance} errs in \emph{opposite} directions:
 it understates $\mathbf{1}^{\transpose}V\mathbf{1}$ on the one whose column sums are
 further from constant, and overstates it on the one whose column sums are nearer.
 Neither the size of the discrepancy nor its sign can therefore be read off the column
 sums; both depend on the whole eigenstructure of $\excitation$ and on
 $\stationaryIntensity$.
 The discrepancy does not depend on $\decay$, since both sides scale as $1/\decay$.
\end{remark}

When the covariance is wanted, \eqref{eq.lyapunov} is a linear system in the
$\numEventTypes(\numEventTypes+1)/2$ distinct entries of $V$,
solved once.
